\documentclass[11pt]{article}

\usepackage[margin=1in]{geometry}
\usepackage{enumitem}
\usepackage{amsmath, amssymb}
\usepackage{hyperref}

\setlist[enumerate]{itemsep=0.8em, topsep=0.8em}
\setlist[itemize]{itemsep=0.3em, topsep=0.3em}

\begin{document}

\begin{center}
  {\Large \textbf{MED 520 Examination}}\\[0.25em]
  {\large \textbf{General Pathology and Disease Mechanisms}}\\[0.5em]
\end{center}

\noindent\textbf{Instructions:}
\begin{itemize}
  \item Answer all questions.
  \item For Questions 1--4, choose the best option.
  \item For Questions 5--8, write brief but precise answers.
\end{itemize}

\vspace{0.5em}

\begin{enumerate}[label=\textbf{\arabic*.}]

  \item Which cellular adaptation involves an increase in cell size in response to increased 
  functional demand?
  \begin{enumerate}[label=(\Alph*)]
    \item Hyperplasia
    \item Hypertrophy
    \item Metaplasia
    \item Dysplasia
  \end{enumerate}

  \item Which type of necrosis is characteristically seen in tuberculosis infection?
  \begin{enumerate}[label=(\Alph*)]
    \item Liquefactive necrosis
    \item Coagulative necrosis
    \item Caseous necrosis
    \item Fat necrosis
  \end{enumerate}

  \item In the inflammatory response, which cells are typically the first to arrive at the site of 
  acute inflammation?
  \begin{enumerate}[label=(\Alph*)]
    \item Lymphocytes
    \item Macrophages
    \item Neutrophils
    \item Eosinophils
  \end{enumerate}

  \item Which tumor marker is most commonly elevated in hepatocellular carcinoma?
  \begin{enumerate}[label=(\Alph*)]
    \item CA-125
    \item PSA (Prostate-Specific Antigen)
    \item AFP (Alpha-Fetoprotein)
    \item CEA (Carcinoembryonic Antigen)
  \end{enumerate}

  \item Explain the difference between benign and malignant tumors in terms of growth patterns, 
  invasion, metastasis, and differentiation. Why is metastatic potential the most critical 
  distinguishing feature?

  \item Describe the pathophysiology of atherosclerosis. What are the major risk factors, and how 
  do they contribute to the formation of atherosclerotic plaques? Explain the role of endothelial 
  dysfunction, lipid accumulation, and inflammation.

  \item Discuss the cellular and molecular mechanisms of apoptosis (programmed cell death). How does 
  it differ from necrosis, and what are the two major pathways that trigger apoptosis? Explain the 
  role of caspases and the significance of apoptosis in development and disease.

  \item Explain the concept of carcinogenesis and the multi-step model of cancer development. Describe 
  the roles of oncogenes, tumor suppressor genes, and genes regulating apoptosis in the transformation 
  from normal cells to malignant cells.

\end{enumerate}

\end{document}
